% Source: 17_Chapter_9_Laying_Her_Sceptre_Down_The_Conversion_o.txt
\bookchapter[Laying Her Sceptre Down]{CHAPTER 9}{LAYING HER SCEPTRE DOWN}{The Conversion of the Empire}

\leadin{U}{NLESS HE HAS SOME} special interest in history or in sculpture, today’s tourist in Rome
will marvel at the size of the Triumphal Arch of Constantine. But if he stops to read his guidebook
carefully he will recognize the Latin at the top of the arch is praising “The Emperor Caesar Flavius
Constantine the Great” for some impressive victory over “the enemy and his whole troop.”

A frieze on the town side of the arch shows the enemy, a rival Caesar named Maxentius, and his
troops falling from the Milvian Bridge into the Tiber River. They are the battered and bleeding
victims of the victorious Constantine.

Eusebius of Caesarea, the Christian historian and biographer of Constantine, compared Maxentius’s
destruction in the Tiber to the Pharaoh’s in the Red Sea. Many contemporaries recognized that
Constantine’s conversion to Christianity and his victory over his enemies was “inspired by the
Godhead” and marked a turning point in history.

Constantine represents the passing of the Age of Catholic Christianity and the beginning of the Age
of the Christian Empire (312–590). Courageous martyrs were a thing of the past. The Christianization
of the empire and the imperial interference in the affairs of the church begins. We can detect the
fallout of these two developments to this day.

How could such a turn of events even come about? Why did the despised and persecuted “superstition”
called Christianity rather suddenly arise from the shadows of Roman society and assume, almost
overnight, the spiritual leadership of the vast and powerful empire?

\section{TURNING AROUND AN EMPIRE}

To catch the significance of this shift in Christianity’s fortunes we need to go back to
Constantine’s predecessor in the imperial palace, a general named Diocletian.

All the signs of a crumbling empire were there in 284, when Diocletian first donned the purple.
Thirty emperors had claimed the throne in the third century, and many others tried. The Roman Senate
scarcely pretended any longer to go through the motions of electing the Caesars. Kinship to a
reigning emperor counted for nothing in determining the succession. On the contrary, the sons or
near relatives of one Caesar usually found themselves in danger of a speedy execution when the next
emperor assumed the crown.

As Hutchinson and Garrison describe the crisis,

\begin{blockquote}
Chaos and anarchy spread throughout the empire. The slaying of one Caesar was a signal to Roman
troops somewhere to acclaim a new ruler. Sometimes the Praetorian Guards stationed in Rome itself
made the choice; sometimes it came from the armies on the frontiers. As the third century drew
toward its close, most thoughtful Romans were in despair. They saw the empire on a swift slide into
ruin and the once proud civilization about to plunge into a barbarian sea.
\end{blockquote}

That of course was what eventually happened. But the collapse did not come in AD 300. In the west it
was delayed for another 175 years. And in the east an empire that claimed to be Roman continued
almost until the day Columbus hauled up anchor to sail to the discovery of the New World. “This
sudden about-face in Rome’s history, from chaos and encroaching disintegration to a new surge of
vigor and stability, was largely the product of one reign, that of Diocletian,” who occupied the
throne for 20 years (284–305), just before Constantine.

Diocletian has not enjoyed what is called a “good press” in Christian circles because he was the
most savage of the persecutors of the church. Given the anarchy he inherited and the revived empire
he passed on to his successors, “Diocletian deserves to be ranked among the truly great emperors.”

The son of slaves in Dalmatia (today’s Balkan area), Diocletian embarked on a military career and
became commander of the army before he was forty. Raised to the purple “by the election of generals
and officers,” he settled any possible competition for the imperial power by leaping on his nearest
rival as he stood before the tribunal of the Senate and running him through. “From that moment on,
this rough-handed soldier proceeded to take the disintegrating empire as it were by the scruff of
its neck and shake it into new life.” He not only turned back the retreat in Germany and along the
Danube; he even reconquered most of distant Britain and Persia.

Diocletian, however, was more than a victorious general. In the long run, his skill as a statesman
proved more significant. He was convinced that the empire was unmanageable. The constant frontier
attacks called for reorganization. So, he divided the imperial power with three others and
established four imperial courts, none of them in Rome. His own court he placed in Nicomedia, in the
northwest corner of Asia Minor. From there he could keep a close watch on the ever-threatening
invaders along the eastern borders.

Diocletian’s plan aimed to protect the empire from the anarchy created by the constant assassination
of emperors. The shrewd old emperor believed that the division of the imperial rule among four
men—two “Augusti,” each with his slightly subordinate “Caesar”—would control the ambitions of his
most likely rivals. Since each had hopes of becoming emperor he would no longer be tempted to
promote his fortunes by killing the older rulers.

\section{CHRISTIANS TO THE FLAMES}

No one seems to know exactly why, but Diocletian, two years before the end of his highly effective
reign, suddenly ordered the most vicious of all persecutions of the Christians. For eighteen years
Diocletian, although himself a convinced and practicing pagan, paid no attention to the growing
Christian power. His court was full of Christian officials, and his wife Prisca and his daughter
Valeria were considered Christians. Impressive church buildings appeared in the principal cities of
the empire, the largest in his capital of Nicomedia.

Then, suddenly, the old emperor ordered his army purged of Christians. Imperial edicts followed,
commanding officials to destroy church buildings, prohibit Christian worship, and burn the
Scriptures. Bishops were rounded up wholesale, imprisoned, tortured, and many put to death, while
the power of the imperial throne was turned loose to wipe out the rest of the Christian community in
blood.

In 305 Diocletian, following a long-established plan, abdicated and forced his fellow Augustus,
Maximian, to do likewise. What history recalls as “the Diocletian persecution” was still raging. In
fact, the new Augustus in the east, Galerius, was more intent than ever on pushing ahead to the
complete extermination of Christianity. Christians said he was the original instigator of the purge.

The pagans themselves, however, were sickened by so much bloodshed. The other new Augustus,
Constantius Chlorus, in far-off Britain, who had never pushed the persecution very hard in his
district of Gaul, suspended all measures against the Christians and began to show them signs of
favor.

In 311, on his deathbed, Galerius realized that his attempt to do away with the upstart religion had
failed. “Thousands upon thousands of terrified Christians had, to be sure, recanted, but other
thousands had stood fast, sealing their faith with their blood.” So eager, in fact, did many
Christians prove to suffer for their faith the bishop of Carthage demanded that those who needlessly
rushed into martyrdom should not be revered as martyrs.

The effect on public opinion throughout the empire was tremendous. “Even the throne could no longer
take the risk of continuing the torturing, maiming, and killing. So, in his last official act,
Galerius, reluctantly, grudgingly, issued an edict of toleration,” and for all practical purposes,
the last and worst persecution of Christians by Rome came to an end.

Upon the death of Galerius, a struggle for imperial power broke loose. In the spring of 312,
Constantine, the son of Constantius Chlorus, advanced across the Alps to dislodge his rival
Maxentius from Italy and to capture Rome. It was a daring gamble; and when he came upon his
militarily superior enemy at the Milvian Bridge, just outside the walls of Rome, he found help in
the God of the Christians. In a dream he saw a cross in the sky and the words, “In this sign
conquer.” This convinced him to advance. When on October 28, 312, he achieved his brilliant victory
over the troops of Maxentius, Constantine looked upon his success as proof of the power of Christ
and the superiority of the Christian religion.

Some historians have considered Constantine’s “conversion” a purely political maneuver. Plenty of
paganism remained. He conspired; he murdered; he even retained his title Pontifex Maximus as head of
the state religious cult. But a purely political conversion is hard to maintain in the light of his
public and private actions. From the year 312, he favored Christianity openly. He allowed Christian
ministers to enjoy the same exemption from taxes as the pagan priests; he abolished executions by
crucifixion; he called a halt to the battles of gladiators as a punishment for crimes; and in 321 he
made Sunday a public holiday. Thanks to his generosity, magnificent church buildings arose as
evidence of his support of Christianity.

\begin{sidebar}{Care for the Dying}
Unbelieving Romans, like many ancients, would abandon unwanted children to die. Girls, who were less
economically desirable, were especially likely victims. Christians holding each person’s worth as
one created in the image of God would rescue these abandoned children. Romans objected to Christians
engaging in this practice, and eventually laws were crafted to penalize believers who rescued the
children. Constantine ruled that the abandoned children belonged to the finder. This ruling, like
many others, favored Christians.
\end{sidebar}

This public Christianity was matched by changes in Constantine’s private life. Making no secret of
his Christian convictions, he had his sons and daughters brought up as Christians and led a
Christian family life. Bishop Eusebius of Nicomedia baptized him shortly before he died in 337.
After his baptism, Constantine refused to wear again the imperial purple and thus left this life
dressed in his white baptismal robes.

The emperor’s conversion inspired Matthew Arnold, the nineteenth-century poet, to write:

\begin{poem}
She heard it, the victorious West,\\
In crown and sword array’d!\\
She felt the void which mined her breast,\\
She shiver’d and obey’d.\\
She veil’d her eagles, snapp’d her sword,\\
And laid her sceptre down;\\
Her stately purple she abhorr’d\\
And her imperial crown.
\end{poem}

\section{A NEW CAPITAL FOR CHRISTIANITY}

Along with his new religion, Constantine provided the Roman world with a new capital city. The
empire’s enemies tended to gather in the east and now the official religion was eastern. A move to
the east was natural. The site Constantine chose could hardly have been more perfect for a city that
would grow rich through trade. The narrow neck of the Bosporus, the waterway between the Black Sea
and the Mediterranean, is a natural crossroads between Asia and Europe, either by land or by water.
Neatly placed to one side of the straits is an ideal natural harbor, narrow enough for a chain to be
stretched across it in the old days to keep out intruders. For centuries a town had been there.
Byzantium it had been called, until the day in 330 when Constantine was ready to inaugurate his new
city. Exactly 1,600 years later, in 1930, the Turks changed its name to Istanbul. But for the
greater part of its long history the place has been known as Constantine’s city, Constantinople, the
capital of the Byzantine (or eastern Roman) Empire.

Today, after all the intervening years, we find it almost impossible to grasp what this change in
imperial leadership meant for the church. Prior to 312, Christianity had been outlawed and
persecuted. Suddenly it was favored and pampered. Constantine thrust it into public life. As a
result, the church re-envisioned its image and mission.

The church historian Eusebius probably spoke for the majority of Christians when he represented the
emperor as the ideal Christian ruler and envisioned the beginning of a new age of salvation. The new
opportunity to preach publicly and to develop unmolested surely meant that God had a new and greater
mission for the church. The divinely ordained moment had arrived for the infusion of public life
with the spirit of Christianity.

Some Christians found special meaning in the fact that Jesus and his message appeared at the very
hour when the Roman Empire provided the world with political, economic, and cultural unity. The
empire seemed to carry the providential task of preparing the way for Christianity to fulfill its
mission to all men. Now, with Constantine, the conversion of the world seemed near.

The advantages for the church were real enough, but there was a price to pay. Constantine ruled
Christian bishops as he did his civil servants and demanded unconditional obedience to official
pronouncements, even when they interfered with purely church matters. There were also the masses who
now streamed into the officially favored church. Prior to Constantine’s conversion, the church
consisted of convinced believers who were willing to bear the risk of being identified as
Christians. Now many came who were politically ambitious, religiously disinterested, and still
half-rooted in paganism. This threatened to produce not only shallowness and permeation by pagan
superstitions but also the secularization and misuse of religion for political purposes.

\section{THE CHURCH COMES TO POWER}

By 380, rewards for Christians had given way to penalties for non-Christians. In that year the
emperor Theodosius made belief in Christianity a matter of imperial command:

\begin{blockquote}
It is Our Will that all the peoples we rule shall practice that religion which the divine Peter the
Apostle transmitted to the Romans. We shall believe in the single Deity of the Father, the Son, and
the Holy Spirit, under the concept of equal majesty and of the Holy Trinity. We command that those
persons who follow this rule shall embrace the name of Catholic Christians. The rest, however, whom
We adjudge demented and insane, shall sustain the infamy of heretical dogmas, their meeting places
shall not receive the name of churches, and they shall be smitten first by divine vengeance and
secondly by the retribution of Our own initiative, which We shall assume in accordance with divine
judgment.
\end{blockquote}

Theodosius takes for granted the close link between his own will and God’s. It was a connection
implicit in the Christian empire.

Church buildings in the Christian empire were carefully designed to emphasize the new hierarchy of
Christ and emperor. The style was borrowed from the East. A Greek traveler visited Persia in the
second century and described a palace in which there was “a hall covered with a dome; the inside was
adorned with sapphires sparkling with a celestial blue brilliance, and standing out against the blue
background of the stones were golden images of the gods, glittering like stars in the firmament.”
This became the pattern for the mosaic-encrusted interiors of Byzantine churches, displaying, if not
“golden images of the gods,” at least God and the “demigod” emperor who represented him on earth.

In the west, farther from the imperial courts, some churchmen dared to challenge the demigod.
Ambrose, Bishop of Milan, was one. The incident that led to the clash was, on the surface, hardly a
matter for imperial concern.

In the year 390 a charioteer in a Greek city was accused of homosexual practices. The governor of
the area threw him into prison but did not count on the reaction of the people. With the chariot
races about to begin, the people asked for the charioteer’s freedom. The governor refused. So the
people rose in arms, murdered the governor, and freed their hero.

Theodosius, then in Milan, was incensed. He ordered that the people be punished, so at another
chariot race in the circus at Thessalonica the gates were closed and the soldiers of Theodosius were
stationed at the entrances. At a signal they fell upon the people. In three hours seven thousand
Thessalonians fell by the sword.

A cry of horror rose through the empire. Ambrose, who regarded himself as the imperial conscience,
felt utterly ashamed. In the name of common humanity and the church, he had to speak out. Theodosius
would have to admit he had committed a crime and repent. Ambrose decided to write Theodosius.

“I cannot deny that you have a zeal for the faith,” he wrote,

\begin{blockquote}
and that you fear God, but you have a naturally passionate spirit which becomes un-governable when
you are excited. I call on you to repent. You can only atone for your sin by tears, by penitence, by
humbling your soul before God. You are a man, and as you have sinned as a man so you must repent. No
angel, no archangel can forgive you. God alone can forgive you, and He forgives only those who
repent.
\end{blockquote}

Ambrose refused the emperor Communion until he had confessed his sin. For a while Theodosius stayed
away from church, but in the end he accepted Ambrose’s terms. In front of a crowded congregation he
took off his splendid imperial robes and asked pardon for his sins. He had to do so on several
occasions until at last, on Christmas Day, Ambrose gave him the sacrament.

It required unusual courage to humiliate a Byzantine emperor. Ambrose had hit upon the weapon—the
threat of excommunication—that the Western church would soon use again and again to humble princes.
But at the center of the Christian empire, in Constantinople, no bishop ever stepped so far out of
line.

Today, as Bamber Gascoigne points out,

\begin{blockquote}
In the Milan church named after St. Ambrose, the services are Roman Catholic—recognizably different
from the form of worship associated with the Byzantine emperors, which we now know as Greek
Orthodox. But orthodox merely means correct; catholic is a word for universal. We might equally well
refer to them as Greek Catholic and Roman Orthodox.
\end{blockquote}

\noindent It was just a case of each side, East and West, claiming to have the right form of
Christianity. In their contrasting attitudes toward the Christian emperors, however, we have a
symbol of their diverging destinies.

\readinghead{Suggestions for Further Reading}
\begin{readinglist}
\entry Baynes, Norman H. Constantine the Great and the Christian Church. 2nd ed. London: Oxford
  University Press, 1972.
\entry Cochrane, Charles Norris. Christianity and Classical Culture. New York: Oxford, 1957.
\entry Coleman, Christopher Bush. Constantine the Great and Christianity. New York: AMS Press, 1968.
\entry Davidson, Ivor. A Public Faith: From Constantine to the Medieval World, AD 312–600, Baker
  History of the Church. Vol. 2. Grand Rapids: Baker, 2005.
\entry Kee, Allistair. Constantine Versus Christ: The Triumph of Ideology. London: SCM Press, 1982.
\entry *Leithart, Peter. Defending Constantine: The Twilight of an Empire and the Dawn of
  Christendom. Downers Grove, IL: InterVarsity Press, 2010.
\entry *Yoder, John Howard. Christian Attitudes to War, Peace, and Revolution. Ed. By Theodore
  Koontz and Andy Alexis-Baker. Grand Rapids: Brazos, 2009.
\end{readinglist}
